% ============================================================
% Section 13: 21st Century & AI/ML/LLMs (Slides 92-99)
% ============================================================
\section{21st Century \& Artificial Intelligence}

% ----------------------------------------------------------
% Slide 92: Section Divider
% ----------------------------------------------------------
{
\setbeamercolor{background canvas}{bg=aiViolet}
\begin{frame}[plain]
  \centering
  \vspace{2cm}
  \textcolor{white}{\Huge\bfseries Mathematics in the}\\[0.4em]
  \textcolor{white}{\Huge\bfseries Age of Machines}\\[1em]
  \textcolor{white!80}{\large 2000--2026 CE}\\[1.5em]
  \textcolor{white!60}{\normalsize Perelman \(\cdot\) Tao \(\cdot\) Mirzakhani \(\cdot\) Uhlenbeck}\\[0.3em]
  \textcolor{white!60}{\normalsize Machine Learning \(\cdot\) Large Language Models}

  \note{
    Transition: ``Turing's theoretical machine became real. And now, machines are
    beginning to DO mathematics.''
    This section covers living mathematicians and the intersection of mathematics
    with AI --- where the 5,000-year story meets the present day.
  }
\end{frame}
}

% ----------------------------------------------------------
% Slide 93: Grigori Perelman and the Poincare Conjecture
% ----------------------------------------------------------
\begin{frame}[shrink=17]{Grigori Perelman --- Solving the Unsolvable, Refusing the Prize}

  \begin{columns}[T]
    \begin{column}{0.5\textwidth}
      \begin{personbox}[Grigori Yakovlevich Perelman (b.\,1966)]{aiViolet}
        \textbf{Origin:} Leningrad (now St.\ Petersburg), Russia\\[3pt]
        \textbf{Key work:} Proved the \alert{Poincar\'e conjecture}
        (2002--2003, published on arXiv), one of the seven Millennium
        Prize Problems.\\[4pt]
        Then \textbf{refused} the \$1\,million prize AND the Fields Medal
        --- the only person to have declined a Fields Medal (as of 2026).\\[4pt]
        He essentially disappeared from public life.
      \end{personbox}

      \vspace{0.3em}
      \begin{insightbox}
        He solved one of the hardest problems in mathematics, turned down
        a million dollars and the highest honour in math, and vanished.
      \end{insightbox}
    \end{column}

    \begin{column}{0.47\textwidth}
      % Poincar\'e conjecture illustrated
      \visualplaceholder[5cm]{Sphere -- rubber band contracts}
    \end{column}
  \end{columns}

  \note{
    ``He solved one of the hardest problems in mathematics, turned down a million
    dollars and the highest honour in math, and essentially disappeared from
    public life.''
    Perelman posted his proof on arXiv (a free preprint server) rather than
    submitting to a journal. He felt the mathematical community's priorities
    were misplaced.
    \verifytag{``Only person ever to decline Fields Medal'' --- add qualifier
    ``as of 2026'' since this could change.}
  }
\end{frame}

% ----------------------------------------------------------
% Slide 94: Terence Tao -- The Mozart of Math
% ----------------------------------------------------------
\begin{frame}[shrink=13]{Terence Tao --- The Mozart of Math}

  \begin{columns}[T]
    \begin{column}{0.5\textwidth}
      \begin{personbox}[Terence Chi-Shen Tao (b.\,1975)]{aiViolet}
        \textbf{Origin:} Adelaide, Australia (Chinese--Hong Kong descent);
        works at UCLA\\[3pt]
        \textbf{Key work:}
        \begin{itemize}\setlength\itemsep{2pt}
          \item Fields Medal (2006)
          \item Work spans: harmonic analysis, PDE, analytic number theory,
                random matrix theory, combinatorics
          \item \alert{Green--Tao theorem} (arXiv 2004, published in
                \textit{Annals of Mathematics} 2008): there are arbitrarily
                long arithmetic progressions in the primes
          \item Pioneer of collaborative, open mathematics (blog, Polymath)
        \end{itemize}
      \end{personbox}
    \end{column}

    \begin{column}{0.47\textwidth}
      % Green--Tao theorem illustrated
      \visualplaceholder[5cm]{Number line with primes highlighted}
    \end{column}
  \end{columns}

  \note{
    Tao was a child prodigy: he scored 760 on the maths SAT at age 8.
    But the emphasis here is on the breadth and depth of his adult work.
    The Green--Tao theorem was first posted to arXiv in 2004 and published
    in the Annals of Mathematics in 2008.
    \verifytag{Green--Tao: arXiv 2004, published Annals 2008.}
    Tao also runs an influential mathematical blog and has been a leader
    in the ``Polymath'' collaborative mathematics movement.
  }
\end{frame}

% ----------------------------------------------------------
% Slide 95: Maryam Mirzakhani -- Breaking Barriers
% ----------------------------------------------------------
\begin{frame}[shrink=25]{Maryam Mirzakhani --- Breaking Barriers}

  \begin{columns}[T]
    \begin{column}{0.52\textwidth}
      \begin{personbox}[Maryam Mirzakhani (1977--2017)]{aiViolet}
        \textbf{Origin:} Tehran, Iran; worked at Stanford University\\[3pt]
        \textbf{Key work:}
        \begin{itemize}\setlength\itemsep{2pt}
          \item \alert{First woman to win the Fields Medal} (2014)
          \item Dynamics and geometry of Riemann surfaces and their
                moduli spaces
          \item Connected hyperbolic geometry, complex analysis, and
                dynamical systems
        \end{itemize}
      \end{personbox}

      \vspace{0.3em}
      \begin{insightbox}
        She died of cancer at 40. Iran honoured her with a postage stamp
        --- showing her without a headscarf.
      \end{insightbox}

      \vspace{0.2em}
      {\small\itshape Her daughter called her work ``painting'' --- she
      would draw huge diagrams on the floor, filling sheets of paper
      with ideas.}
    \end{column}

    \begin{column}{0.45\textwidth}
      \visualplaceholder[5cm]{Maryam Mirzakhani working on large sheets\\
        of paper on the floor\\(Stanford / photographer credit needed)}

      \vspace{0.5em}
      % Riemann surface visualization
      \visualplaceholder[5cm]{Genus-0 (sphere)}
    \end{column}
  \end{columns}

  \note{
    ``She was the first woman in history to win the Fields Medal. She died of cancer
    at 40, and Iran put her picture on a postage stamp --- without a headscarf.''
    Mirzakhani's work is deep and abstract, but her approach was visual and intuitive.
    She once described her mathematical work as being ``like a novel with different
    characters, and you are getting to know them better.''
    Her loss at age 40 was mourned across the global mathematical community.

    CALLBACK -- WOMEN'S THREAD COMPLETED: ``Remember what we said at Noether's
    slide? Hypatia, Kovalevskaya, Noether -- denied, blocked, killed. Mirzakhani
    broke through. First woman to win the Fields Medal. And she died at 40. The
    fight is not over -- but the barrier is broken.''
  }
\end{frame}

% ----------------------------------------------------------
% Slide 96: Karen Uhlenbeck -- Geometry and Gauge Theory
% ----------------------------------------------------------
\begin{frame}[shrink=13]{Karen Uhlenbeck --- Geometry and Gauge Theory}

  \begin{columns}[T]
    \begin{column}{0.52\textwidth}
      \begin{personbox}[Karen Keskulla Uhlenbeck (b.\,1942)]{aiViolet}
        \textbf{Origin:} Cleveland, Ohio, USA; worked at University of Texas
        at Austin and Princeton\\[3pt]
        \textbf{Key work:}
        \begin{itemize}\setlength\itemsep{2pt}
          \item \alert{First woman to win the Abel Prize} (2019)
          \item ``Pioneering achievements in geometric partial differential
                equations, gauge theory, and integrable systems''
          \item Work underpins modern mathematical physics, including the
                mathematics behind Yang--Mills gauge theories
          \item Co-founded the Women and Mathematics program at the
                Institute for Advanced Study
        \end{itemize}
      \end{personbox}
    \end{column}

    \begin{column}{0.45\textwidth}
      \visualplaceholder[4.5cm]{Portrait of Karen Uhlenbeck\\(Wikimedia Commons)}

      \vspace{0.5em}
      % Gauge theory visualization
      \visualplaceholder[5cm]{A surface with connections (arrows showing parallel transpor}
    \end{column}
  \end{columns}

  \note{
    ``She won the Abel Prize --- one of math's highest honours --- for work that
    connects geometry to the fundamental forces of nature. She also spent decades
    fighting to make mathematics more welcoming to women.''
    Uhlenbeck's work is highly technical, but the key idea: she proved foundational
    results about the mathematical structures (gauge fields) that physicists use
    to describe electromagnetism, the strong force, and the weak force.
    Together with Mirzakhani (Fields Medal) and Noether, she represents the
    growing recognition of women's contributions to mathematics.
  }
\end{frame}

% ----------------------------------------------------------
% Slide 97: Machine Learning and the New Mathematics
% ----------------------------------------------------------
\begin{frame}[shrink=25]{Machine Learning --- The New Mathematics}

  \begin{columns}[T]
    \begin{column}{0.42\textwidth}
      {\small The mathematics of machine learning:}

      \vspace{0.3em}
      \begin{tcolorbox}[colback=aiViolet!5, colframe=aiViolet, rounded corners,
        boxrule=0.7pt, left=3pt, right=3pt, top=2pt, bottom=2pt, fonttitle=\bfseries\small,
        title={Math Inside Every AI}]
        \small
        \begin{itemize}\setlength\itemsep{2pt}
          \item \textbf{Linear algebra} --- neural networks are matrix operations
          \item \textbf{Calculus} --- gradient descent = optimisation via derivatives
          \item \textbf{Probability} --- Bayesian inference, statistical learning
          \item \textbf{Information theory} --- Shannon's legacy
        \end{itemize}
      \end{tcolorbox}

      \vspace{0.3em}
      {\scriptsize\textbf{Key figures:}\\
      \textbf{Geoffrey Hinton} (b.\,1947, UK/Canada) --- deep learning, backpropagation.
      Nobel Prize in Physics 2024 (shared with John Hopfield).\\[2pt]
      \textbf{Yoshua Bengio} (b.\,1964, France/Canada) --- deep learning.
      Turing Award 2018 (shared).\\[2pt]
      \textbf{Yann LeCun} (b.\,1960, France/USA) --- convolutional neural
      networks. Turing Award 2018 (shared).}
    \end{column}

    \begin{column}{0.55\textwidth}
      % TikZ: Neural network annotated with math operations
      \visualplaceholder[5cm]{Input layer}
    \end{column}
  \end{columns}

  \note{
    ``Every AI you interact with --- every chatbot, every image generator --- runs
    on the mathematics we have been studying today: algebra, calculus, linear algebra,
    probability. 5,000 years of math, all in one machine.''
    \verifytag{Hinton Nobel Prize Physics 2024 shared with Hopfield. Verify exact
    citation phrasing.}
    \verifytag{Bengio/LeCun Turing Award: awarded 2018, announced March 2019?
    Verify.}
    Timeline: Perceptron (Rosenblatt, 1958) $\to$ backpropagation (Rumelhart,
    Hinton, Williams, 1986) $\to$ deep learning revolution (2012+) $\to$
    Transformers (2017) $\to$ modern LLMs.

    NARRATIVE CLIMAX: Do NOT deliver this as a technology overview. Deliver it
    as the punchline of the entire lecture. ``Look at this neural network
    diagram. Matrix multiplication -- that is linear algebra, which we traced
    from Babylon to Gauss. Gradient descent -- that is calculus, Newton and
    Leibniz. Probability -- that is Pascal and Laplace. Information theory --
    that is Shannon. EVERYTHING we have studied today, 5,000 years of human
    genius, is running inside this machine. And the machine is now discovering
    NEW mathematics that no human has ever seen. The relay has a new kind of
    runner.''
  }
\end{frame}

% ----------------------------------------------------------
% Slide 98: The Story Continues
% ----------------------------------------------------------
\begin{frame}[shrink=25]{The Story Continues}

  \begin{columns}[T]
    \begin{column}{0.5\textwidth}
      {\normalsize\bfseries Mathematics is not finished.}

      \vspace{0.5em}
      \begin{tcolorbox}[colback=aiViolet!5, colframe=aiViolet, rounded corners,
        boxrule=0.7pt, left=3pt, right=3pt, top=3pt, bottom=3pt,
        fonttitle=\bfseries\small, title={Open Problems}]
        \small
        \begin{itemize}\setlength\itemsep{3pt}
          \item \textbf{Riemann hypothesis} (1859) --- \$1M prize
          \item \textbf{P vs NP} --- can every problem whose answer is
                quickly verifiable also be quickly solved?
          \item \textbf{Navier--Stokes} --- do smooth solutions always exist?
          \item \textbf{Birch--Swinnerton-Dyer} conjecture
        \end{itemize}
      \end{tcolorbox}

      \vspace{0.3em}
      \begin{tcolorbox}[colback=accentOrange!5, colframe=accentOrange, rounded corners,
        boxrule=0.7pt, left=3pt, right=3pt, top=3pt, bottom=3pt,
        fonttitle=\bfseries\small, title={AI Discovers Mathematics}]
        \small
        \begin{itemize}\setlength\itemsep{2pt}
          \item DeepMind's \textbf{AlphaGeometry} --- solves olympiad geometry
          \item \textbf{Automated theorem proving} (Lean, Coq)
          \item AI-assisted conjecture generation
        \end{itemize}
      \end{tcolorbox}
    \end{column}

    \begin{column}{0.47\textwidth}
      % Mega-timeline fully filled in (return to Slide 3 timeline)
      \visualplaceholder[5cm]{Main axis}
    \end{column}
  \end{columns}

  \note{
    Return to the mega-timeline from Slide 3, now fully filled in.
    ``The story of mathematics is the story of human curiosity. It started with
    counting, it gave us civilisation, and it is not over. The next chapter has
    not been written yet.''
    Point out that AI is now being used to DISCOVER new mathematics --- we may
    be entering a fundamentally new era.
    The question mark at the end is for the students: the next breakthrough could
    come from anyone.
  }
\end{frame}

% ----------------------------------------------------------
% Slide 99: Key Themes Recap
% ----------------------------------------------------------
\begin{frame}[shrink=11]{Key Themes --- What Have We Learned?}

  \begin{columns}[T]
    \begin{column}{0.5\textwidth}
      \begin{enumerate}\setlength\itemsep{6pt}
        \item[\textcolor{aiViolet}{\faGlobeAmericas}]
          \textbf{Math has no single homeland} --- built by people in Central
          Africa, Babylon, Egypt, Mesoamerica, Greece, India, China, the
          Islamic world, Europe, and every corner of the globe

        \item[\textcolor{aiViolet}{\faArrowsAltH}]
          \textbf{Mathematical ideas travel} --- the same concept was often
          discovered independently in multiple civilisations

        \item[\textcolor{aiViolet}{\faUser}]
          \textbf{People matter} --- behind every theorem is a human story

        \item[\textcolor{aiViolet}{\faFistRaised}]
          \textbf{Barriers have been broken} --- from Hypatia to Noether to
          Kovalevskaya to Mirzakhani to Uhlenbeck, women have fought for
          and won their place in mathematics
      \end{enumerate}
    \end{column}

    \begin{column}{0.47\textwidth}
      % World map with all mathematicians plotted
      \visualplaceholder[5cm]{Simplified world outline}

      \vspace{0.2em}
      \begin{center}
        {\scriptsize\itshape Mathematics was never the achievement of one culture.\\
        It is humanity's collective masterpiece.}
      \end{center}
    \end{column}
  \end{columns}

  \note{
    The world map shows all 84 mathematicians mentioned in the lecture, plotted
    by their place of origin.
    Emphasise the global nature of mathematics --- no single civilisation has a
    monopoly on mathematical genius.
    Five women highlighted: Hypatia, Noether, Kovalevskaya, Mirzakhani, Uhlenbeck.
    ``Five thousand years. Dozens of civilisations. Hundreds of brilliant minds.
    And the story is not over.''
  }
\end{frame}
